\renewcommand{\theposttitle}{Why failing is not what you think}
\renewcommand{\thepostdate}{2025-04-03}
\vspace*{\fill}
\begin{center}
\begin{minipage}{0.92\textwidth}
\small
Engaging in creativity is like building a LEGO set without instructions.

With each new piece, the whole thing warps and twists. You have a vague sense of the shape, but there inevitably is a point where key pieces go missing and you have no idea how to move forward.

The lights turn off and you start running into brick walls that come out of nowhere. This dark void is called \textbf{“The Zone of Ambiguity,”} a temporary but necessary space.
\end{minipage}
\end{center}
\vspace{4mm}
\begin{center}
\includegraphics[width=\textwidth, height=0.55\textheight, keepaspectratio]{C:/Users/James/Documents/GitHub/VisualIdeasBook/data/images/311e497de466902e4e9348ada58d6317.png}
\end{center}
\vspace{2mm}
\begin{center}
\begin{minipage}{0.92\textwidth}
\small
The only difference between someone who considers themselves creative and someone who doesn’t is the confidence that the \textbf{Zone of Ambiguity is a temporary space.}

And you build confidence through seeing you can in fact can find your way through. With each rep, that gap between ambiguity and insight shrinks smaller, but it will always exist.
\end{minipage}
\end{center}
\vspace*{\fill}